\begin{remark}[7.4.8]
\label{IV.2.s7.4.8-fr}
Los resultados anteriores plantean los problemas siguientes:
\begin{enumerate}
\item[A)] Sea $A$ un anillo de Zariski \emph{completo}, y sea
$\mathfrak J$ un ideal de definición de $A$; si el anillo
$A/\mathfrak J$ es un $\boldsymbol P$-anillo, ¿lo es también $A$? De
ello resultaría que, para todo $\boldsymbol P$-anillo noetheriano $A$ y
todo ideal $\mathfrak J$ de $A$, el completado separado $\widehat A$
para la topología $\mathfrak J$-preádica sería también un
$\boldsymbol P$-anillo.
\item[B)] Sea $k$ un cuerpo valuado completo no discreto; se
llama también anillo de \emph{series formales restringidas}
$k\{T_1,\ldots,T_n\}$ al subanillo del anillo de series formales
$k[[T_1,\ldots,T_n]]$ formado por las series cuyos coeficientes
\emph{tienden a} 0. ¿Es tal anillo un $\boldsymbol P$-anillo?
\item[C)] Si $A$ es un $\boldsymbol P$-anillo linealmente
topologizado y $S$ un conjunto multiplicativo de $A$, ¿son los anillos
$A\{S^{-1}\}$ ($\mathbf{0}_{\mathrm I}$, 7.6.1) y $A_{\{S\}}$
($\mathbf{0}_{\mathrm I}$, 7.6.15) $\boldsymbol P$-anillos?
\end{enumerate}
\end{remark}
